% ============================================================================
\clearpage
\section*{v4.9.0-r31 Source-Only: The Markov Cascade as a Test of Self-Stabilisation, and a Status Note to Theorem~\texorpdfstring{\ref{thm:dark-energy}}{dark-energy}}
\addcontentsline{toc}{section}{v4.9.0-r31 Source-Only: The Markov Cascade as a Test of Self-Stabilisation, and a Status Note to the Dark-Energy Theorem}
\label{sec:v490-r31-markov-self-stabilisation}
% ============================================================================

\paragraph{Purpose.}
r30 noted that Theorem~\ref{thm:dark-energy} calls \(\Omega_M\) and \(\Omega_\Lambda\) Markov transition probabilities, while the compensator
composes amplitudes.  On 2026-10-22 the author asked for a status note.  He explained that the Markov chain was meant to predict the
probabilities better and to test whether the cascade can stabilise itself.  This layer runs that test on the chain of the proof and
places the note under the revision rule of r25.  The script is \texttt{llm/k68\_markov\_kaskade\_stabil.py}.  The note is applied by
\texttt{llm/apply\_patches.py} from \texttt{llm/revision\_spec\_v490\_r31.py}, so that the GitHub candidate can receive the same text.
No new prediction is made, and nothing is compared with data.

\noindent\textbf{Source-only status.}
This v4.9.0-r31 layer is source-only.  It does not build a PDF, does not assign a new DOI,
and does not authorize a public release.  The v4.9.0 Version DOI
\href{https://doi.org/10.5281/zenodo.23534954}{10.5281/zenodo.23534954} and the Concept DOI
\href{https://doi.org/10.5281/zenodo.19210728}{10.5281/zenodo.19210728} are unchanged by it.

\subsection*{v4.9.0-r31 the test}
\label{subsec:v490-r31-test}

The proof of Theorem~\ref{thm:dark-energy} uses the chain \(S^8\to S^4\) with probability \(\bsym^2\), \(S^8\to S^2\) with \(1-\bsym^2\),
\(S^4\to S^2\) with \(1\), and \(S^2\) absorbing.

\paragraph{Theoremlet v4.9.0-r31.1.}
\begin{enumerate}[leftmargin=2em]
\item[(a)] As given, the chain does not stabilise.  After three steps every distribution has reached \(S^2\), and the only stationary
distribution is \((p_2,p_4,p_8)=(1,0,0)\), the photon.
\item[(b)] With a return path \(S^2\to S^8\) of rate \(r>0\) the chain is irreducible and has the unique stationary distribution
\[
(p_2,p_4,p_8)=\frac{(1,\ \bsym^2r,\ r)}{1+r+\bsym^2r}.
\]
Every start converges to it.  At \(\bsym^2=1/3\), \(r=1/2\) the other eigenvalues have modulus \(0.41\), so the approach is a damped
oscillation.
\item[(c)] In the stationary state the flux out of \(S^8\) splits into \(\bsym^2\) towards \(S^4\) and \(1-\bsym^2\) towards \(S^2\), for every return
rate \(r\).  The fractions that the theorem names \(\Omega_M\) and \(\Omega_\Lambda\) are therefore stable flux fractions, independent of how
fast the return is.
\end{enumerate}

\paragraph{Proof sketch.}
Stationary equations \(\pi Q=\pi\), \(\sum\pi=1\), solved symbolically; convergence checked from \(S^8\) over \(200\) steps at
\(r=1/10,1/2,1\); eigenvalues \(1\) and \(-\tfrac14\pm i\sqrt{15}/12\) at \(\bsym^2=1/3\), \(r=1/2\).  \(\square\)

\paragraph{Reading.}
The author's test gives a clear answer.  The closed cascade does not stabilise itself; it ends in the photon.  With a return path it
does, and then the matter fraction \(\bsym^2\) is a stable property of the flow.  This is the same pattern on the level of the strata as on
the doublet in r27 and r28.  Closed, everything flows to the photon.  Replenished, a stable state arises.  A natural return path is the
pair creation from the vacuum of r29.  The rate \(r\) remains open, but the fractions do not depend on it.  The Markov chain is a
tool for these probabilities between the strata.  Within one step, amplitudes compose through the compensator (r30).

\subsection*{v4.9.0-r31 the status note}
\label{subsec:v490-r31-note}

The status note under Theorem~\ref{thm:dark-energy} is extended (rule C of r25).  It keeps the earlier content and adds the author's
purpose for the Markov chain, the result r31.1, and the reading of r30.  The text of the theorem is unchanged.

\subsection*{v4.9.0-r31 status in the three frames}
\label{subsec:v490-r31-frames}

\begin{center}
\small
\begin{tabular}{@{}>{\raggedright\arraybackslash}p{0.30\textwidth}>{\raggedright\arraybackslash}p{0.62\textwidth}@{}}
\toprule
result & M / T / P\\
\midrule
r31.1(a) absorption & the absorbing stratum \(S^2\) is the photon ground state of the main body; the chain ends there in every frame description\\
r31.1(b), (c) & \(\bsym^2\) is a matter-frame quantity (r22); the stationary fractions hold in M, and are frame-free between M and \(T|_R\)\\
\bottomrule
\end{tabular}
\end{center}

\subsection*{v4.9.0-r31 register entries}
\label{subsec:v490-r31-register}

\begin{center}
\small
\begin{tabular}{@{}>{\raggedright\arraybackslash}p{0.34\textwidth}>{\raggedright\arraybackslash}p{0.58\textwidth}@{}}
\toprule
statement & status after v4.9.0-r31\\
\midrule
Markov cascade of Thm.~\ref{thm:dark-energy} & author's tool for probabilities and self-stabilisation; closed: no stable state besides the photon; with return path: unique stationary state, flux fractions \(\bsym^2\), \(1-\bsym^2\) for every rate (r31.1)\\
status note to Thm.~\ref{thm:dark-energy} & extended (rule C of r25)\\
return path & proposed: pair creation from the vacuum (r29); rate open, fractions independent of it\\
\bottomrule
\end{tabular}
\end{center}

\subsection*{v4.9.0-r31 ledger and guard update}
\label{subsec:v490-r31-ledger-guard}

\[
R_{\rm Mk}^{v490\text{-}r31}=(0,0,1)
\]
for (self-stabilisation of the cascade tested, status note placed, rate of the return path open).

The layer spends no guard:
\[
R_{\rm glob}^{(v490\text{-}r31,+)}=(0,1,1,1,1),
\qquad
\Gamma_{\rm guard}^{v490\text{-}r31}=(1,1,1,1,1),
\qquad
\beta_4=1/256\ \text{unspent}.
\]
It does not imply exact finite \(C5\) zero closure, exact all-readout terminality,
\(C6/T2\), a global photon-clock/RH theorem, or global all-readout closure.

